\documentclass[a4paper,12pt]{ltjsarticle}

\usepackage{luatexja}
\usepackage{geometry}
\usepackage{hyperref}
\usepackage{listings}

\geometry{
top=25mm,
bottom=25mm,
left=25mm,
right=25mm
}

\lstset{
basicstyle=\ttfamily\small,
breaklines=true,
frame=single,
numbers=left,
numberstyle=\tiny
}

\title{Webサービス制作 最終レポート}

\author{
学生番号：25G1100\\
氏名：生天目 瑠美
}

\date{\today}

\begin{document}

\maketitle

\section{公開URL}

\url{http://52.231.33.164/index.html}

\section{制作目的}

本作品は，一見すると可愛らしいペット紹介サイトでありながら，特定の操作を行うことでホラー演出へと変化するWebサービスの制作を目的とした．

通常画面では犬やインコなどのペット情報を閲覧でき，ペットごとの特徴を追加・削除することができる．また，サイト内には隠された導線を設けており，行方不明画面や呪い画面へ遷移することで，かわいらしい雰囲気から恐怖演出へ変化する仕組みを実装した．

本制作を通して，Node.jsおよびExpressを利用したWebサーバの構築方法や，APIを利用したデータ管理，EJSによる動的な画面生成について学ぶことを目的とした．

\section{サーバー側プログラム}

以下にサーバー側プログラム（pet.js）を示す．

\begin{lstlisting}%ここに pet.js を丸ごと貼り付ける

"use strict";

const express = require("express");
const app = express();
const PORT = 80;

let pets = require("./petData");

// EJSを使うための設定
app.set("view engine", "ejs");

app.use(express.json());
// 画像や音楽、静的ファイルを読み込めるようにする設定
app.use(express.static(__dirname));

// ペット一覧データ（API）
app.get("/api/pets", (req, res) => {
  res.json(pets);
});

// ペット詳細データ（API）
app.get("/api/pets/:id", (req, res) => {
  const pet = pets.find(p => p.id === Number(req.params.id));
  res.json(pet);
});

// 好きなポイント追加（POST）
app.post("/api/pets/:id/likes", (req, res) => {
  const pet = pets.find(p => p.id === Number(req.params.id));
  pet.likes.push(req.body.text);
  res.json(pet);
});

// 好きなポイント削除（DELETE）
app.delete("/api/pets/:id/likes/:index", (req, res) => {
  const pet = pets.find(p => p.id === Number(req.params.id));
  pet.likes.splice(Number(req.params.index), 1);
  res.json(pet);
});


/* ========================================================
   ここから：EJS（画面表示）用のルーティング設定
   これでキレイなURL（/pet/1 などの形）でアクセスできるようになります
======================================================== */

// ① ペット詳細画面 (例: http://IPアドレス:8081/pet/1)
app.get("/pet/:id", (req, res) => {
  const pet = pets.find(p => p.id === Number(req.params.id));
  if (!pet) return res.status(404).send("ペットが見つかりません");
  // views/detail.ejs を表示、同時にpetデータを渡す
  res.render("detail", { pet: pet });
});

// ② 行方不明画面 (例: http://IPアドレス:8081/missing/1)
app.get("/missing/:id", (req, res) => {
  const pet = pets.find(p => p.id === Number(req.params.id));
  if (!pet) return res.status(404).send("ペットが見つかりません");
  // views/missing.ejs を表示
  res.render("missing", { pet: pet });
});

// ③ 呪い画面 (例: http://IPアドレス:8081/curse/1)
app.get("/curse/:id", (req, res) => {
  const pet = pets.find(p => p.id === Number(req.params.id));
  if (!pet) return res.status(404).send("ペットが見つかりません");
  // views/curse.ejs を表示
  res.render("curse", { pet: pet });
});


app.listen(PORT, () => {
  console.log(`http://localhost:${PORT}`);
});

\end{lstlisting}

\end{document}
